% ============================================================================
% preamble.tex — 《抽象代数导论：面向模形式》
% ============================================================================

% --- 中文支持 (XeLaTeX) ---
\usepackage{ctex}
\usepackage{fontspec}

% --- 页面布局 ---
\usepackage[
  a4paper,
  top=2.5cm, bottom=2.5cm,
  left=2.5cm, right=2.5cm
]{geometry}

% --- 数学 ---
\usepackage{amsmath, amssymb, amsthm}
\usepackage{mathrsfs}
\usepackage{mathtools}
\usepackage{bm}
\usepackage{enumitem}

% --- 图表 ---
\usepackage{graphicx}
\usepackage{float}
\usepackage{booktabs}
\usepackage{longtable}
\usepackage{tabularx}

% --- 绘图 ---
\usepackage{tikz}
\usetikzlibrary{
  positioning, calc, arrows.meta,
  decorations.pathreplacing,
  patterns, matrix,
}
\usepackage{tikz-cd}

% --- 交叉引用与超链接 ---
\usepackage{hyperref}
\hypersetup{
  colorlinks=true,
  linkcolor=red!60!black,
  citecolor=green!40!black,
  urlcolor=cyan!50!black,
}
\usepackage[nameinlink]{cleveref}

% --- 索引 ---
\usepackage{makeidx}
\makeindex

% --- 颜色 ---
\usepackage{xcolor}

% --- 彩色定理框 ---
\usepackage[most]{tcolorbox}
\tcbuselibrary{breakable, skins}

% ============================================================================
% 定理环境
% ============================================================================

\theoremstyle{plain}
\newtheorem{theorem}{定理}[chapter]
\newtheorem{proposition}[theorem]{命题}
\newtheorem{lemma}[theorem]{引理}
\newtheorem{corollary}[theorem]{推论}

\theoremstyle{definition}
\newtheorem{definition}[theorem]{定义}
\newtheorem{example}[theorem]{例}

\theoremstyle{remark}
\newtheorem{remark}[theorem]{注}
\newtheorem{notation}[theorem]{记号}

% --- tcolorbox 包装 ---
\tcolorboxenvironment{theorem}{
  enhanced, breakable, arc=3pt,
  before skip=10pt, after skip=10pt,
  colback=red!4, colframe=red!40!black,
  boxrule=0.6pt, left=6pt, right=6pt, top=6pt, bottom=6pt,
}
\tcolorboxenvironment{proposition}{
  enhanced, breakable, arc=3pt,
  before skip=10pt, after skip=10pt,
  colback=red!4, colframe=red!35!black,
  boxrule=0.6pt, left=6pt, right=6pt, top=6pt, bottom=6pt,
}
\tcolorboxenvironment{lemma}{
  enhanced, breakable, arc=3pt,
  before skip=10pt, after skip=10pt,
  colback=red!4, colframe=red!35!black,
  boxrule=0.6pt, left=6pt, right=6pt, top=6pt, bottom=6pt,
}
\tcolorboxenvironment{corollary}{
  enhanced, breakable, arc=3pt,
  before skip=10pt, after skip=10pt,
  colback=red!4, colframe=red!35!black,
  boxrule=0.6pt, left=6pt, right=6pt, top=6pt, bottom=6pt,
}
\tcolorboxenvironment{definition}{
  enhanced, breakable, arc=3pt,
  before skip=10pt, after skip=10pt,
  colback=green!4, colframe=green!40!black,
  boxrule=0.6pt, left=6pt, right=6pt, top=6pt, bottom=6pt,
}
\tcolorboxenvironment{example}{
  enhanced, breakable, arc=3pt,
  before skip=10pt, after skip=10pt,
  colback=orange!4, colframe=orange!50!black,
  boxrule=0.6pt, left=6pt, right=6pt, top=6pt, bottom=6pt,
}
\tcolorboxenvironment{remark}{
  enhanced, breakable, arc=3pt,
  before skip=10pt, after skip=10pt,
  colback=gray!6, colframe=gray!40,
  boxrule=0.6pt, left=6pt, right=6pt, top=6pt, bottom=6pt,
}

% ============================================================================
% 习题环境
% ============================================================================

\newcommand{\basic}{$\bigstar$}
\newcommand{\medium}{$\bigstar\bigstar$}
\newcommand{\hard}{$\bigstar\bigstar\bigstar$}

% ============================================================================
% 常用数学命令
% ============================================================================

% 数域
\newcommand{\N}{\mathbb{N}}
\newcommand{\Z}{\mathbb{Z}}
\newcommand{\Q}{\mathbb{Q}}
\newcommand{\R}{\mathbb{R}}
\newcommand{\C}{\mathbb{C}}
\newcommand{\F}{\mathbb{F}}

% 上半平面
\newcommand{\HH}{\mathbb{H}}

% 矩阵群
\newcommand{\SL}{\mathrm{SL}}
\newcommand{\GL}{\mathrm{GL}}
\newcommand{\PSL}{\mathrm{PSL}}
\newcommand{\PGL}{\mathrm{PGL}}
\newcommand{\Sp}{\mathrm{Sp}}
\newcommand{\Mat}{\mathrm{Mat}}
\newcommand{\Orth}{\mathrm{O}}
\newcommand{\SO}{\mathrm{SO}}
\newcommand{\Unit}{\mathrm{U}}
\newcommand{\SU}{\mathrm{SU}}
\newcommand{\diag}{\mathrm{diag}}

% 同余子群
\newcommand{\GammaO}{\Gamma_0}
\newcommand{\GammaI}{\Gamma_1}

% 模形式空间
\newcommand{\Mk}{\mathcal{M}}
\newcommand{\Sk}{\mathcal{S}}

% 代数结构
\newcommand{\Aut}{\mathrm{Aut}}
\newcommand{\Gal}{\mathrm{Gal}}
\newcommand{\Frob}{\mathrm{Frob}}
\newcommand{\Hom}{\mathrm{Hom}}
\newcommand{\End}{\mathrm{End}}
\newcommand{\Tr}{\mathrm{Tr}}
\newcommand{\Nm}{\mathrm{Nm}}
\providecommand{\Ann}{\mathrm{Ann}}
\newcommand{\ord}{\mathrm{ord}}
\newcommand{\Id}{\mathrm{Id}}
\newcommand{\Ker}{\mathrm{Ker}}
\newcommand{\Ima}{\mathrm{Im}}
\newcommand{\Char}{\mathrm{char}}
\newcommand{\Spec}{\mathrm{Spec}}

\DeclareMathOperator{\im}{im}
\DeclareMathOperator{\coker}{coker}
\DeclareMathOperator{\rank}{rank}
\DeclareMathOperator{\nullity}{nullity}
\DeclareMathOperator{\sgn}{sgn}
\DeclareMathOperator{\lcm}{lcm}
\DeclareMathOperator{\Stab}{Stab}
\DeclareMathOperator{\Orb}{Orb}
\DeclareMathOperator{\Fix}{Fix}
\DeclareMathOperator{\Syl}{Syl}
\DeclareMathOperator{\Rad}{Rad}

% 内积
\newcommand{\inner}[2]{\langle #1, #2 \rangle}

% 矩阵简写
\newcommand{\mat}[1]{\begin{pmatrix} #1 \end{pmatrix}}
\newcommand{\dmat}[1]{\begin{vmatrix} #1 \end{vmatrix}}

% 正合列
\newcommand{\ses}[3]{0 \to #1 \to #2 \to #3 \to 0}
